
\documentclass[addpoints]{exam}
\usepackage{amsmath}
\usepackage{amssymb}
\usepackage{color}
\usepackage{siunitx}
\usepackage{booktabs}

% \printanswers

\newcommand{\event}{Final Exam}
\newcommand{\tournament}{}
\def\type{0}

\setlength\fillinlinelength{\textwidth}
\CorrectChoiceEmphasis{\color{red}\bfseries}
\SolutionEmphasis{\color{red}\bfseries}

\setlength\answerclearance{2pt}
\pagestyle{headandfoot}
\runningheadrule
\runningheader{\event}{\tournament}{\ifnum\type=1 Class Set Test \else Full Name:\hspace{1cm} \fi}
\runningfootrule
\runningfooter{}{Page \thepage\ of \numpages}{}

\begin{document}
\begin{center}
    {\huge Grade 12 Foundations for College Math \\ \event
    \ifprintanswers \space Key
    \else \space \ifnum\type<2 Test \fi \ifnum\type=1 \textbf{(CLASS SET)} \fi
          \ifnum\type=2 Answer Sheet \fi
    \fi \\ \vspace{3mm}\tournament\vspace{1mm}}
\end{center}

\vspace{.25\textheight}

\begin{center}
    First Name: \ifprintanswers
    \underline{\hspace{3.5cm}}\textbf{KEY}\underline{\hspace{5.5cm}}\\\vspace{5mm}
    \else \underline{\hspace{10cm}}\\ \vspace{5mm} \fi
    Last Name: \underline{\hspace{10cm}}\\ \vspace{5mm}
\end{center}

\noindent\textbf{\underline{Directions:}}
\begin{itemize}
    \ifnum\type=1 \item \textbf{This test is a class set.} Do not write on this paper. \fi
    \item Show all work for full marks. Round dollar amounts to the nearest cent.
    \item The exam is designed to be completed in 75 minutes.
\end{itemize}
\begin{center}
\textit{For grading use only}\\
\multirowgradetable{1}[pages]
\end{center}

\newpage
\begin{questions}

\section*{Part A --- Multiple Choice (20 marks)}
\vspace{5mm}

%--- Financial Mathematics (4) ---
\question[1] The future value of an annuity formula $FV = PMT \cdot \dfrac{(1+i)^n - 1}{i}$ is used when:
\begin{choices}
    \choice Finding the present lump sum needed now
    \correctchoice Finding the accumulated value of regular deposits
    \choice Calculating simple interest
    \choice Finding the monthly payment on a loan
\end{choices}

\question[1] \$500 is deposited monthly into an account at 6\% per year compounded monthly
for 10 years. The number of compounding periods is:
\begin{choices}
    \choice $10$
    \choice $60$
    \correctchoice $120$
    \choice $72$
\end{choices}

\question[1] Which formula gives the monthly mortgage payment?
\begin{choices}
    \choice $FV = PMT \cdot \dfrac{(1+i)^n - 1}{i}$
    \correctchoice $PV = PMT \cdot \dfrac{1-(1+i)^{-n}}{i}$, solved for $PMT$
    \choice $A = P(1+i)^n$
    \choice $I = Prt$
\end{choices}

\question[1] A \$200\,000 mortgage at 5\% per year compounded semi-annually, 25 years,
monthly payments — the effective monthly rate is:
\begin{choices}
    \choice $\dfrac{0.05}{12}$
    \correctchoice $(1 + 0.025)^{1/6} - 1$
    \choice $\dfrac{0.05}{2}$
    \choice $0.05 \div 25$
\end{choices}

%--- Modelling (4) ---
\question[1] A data set is best modelled by a quadratic function when:
\begin{choices}
    \choice The data increases by a constant amount each period
    \correctchoice The data increases then decreases (or vice versa), forming a curve
    \choice The data doubles each period
    \choice The data oscillates regularly
\end{choices}

\question[1] The correlation coefficient $r = -0.88$ indicates:
\begin{choices}
    \choice Weak positive correlation
    \choice No correlation
    \correctchoice Strong negative correlation
    \choice Weak negative correlation
\end{choices}

\question[1] For the exponential model $y = 3(2)^x$, the doubling time is:
\begin{choices}
    \choice $2$
    \correctchoice $1$
    \choice $3$
    \choice $0.5$
\end{choices}

\question[1] The sinusoidal model $y = 4\sin(2x) + 3$ has amplitude:
\begin{choices}
    \choice $2$
    \correctchoice $4$
    \choice $3$
    \choice $7$
\end{choices}

%--- Trigonometry and Geometry (4) ---
\question[1] The Cosine Law is used when you know:
\begin{choices}
    \choice Two angles and a side
    \choice One side and all angles
    \correctchoice Two sides and the included angle (SAS) or all three sides (SSS)
    \choice Two sides and a non-included angle
\end{choices}

\question[1] A cylinder with $r = 2$ m and $h = 5$ m has total surface area:
\begin{choices}
    \choice $14\pi$ m$^2$
    \choice $20\pi$ m$^2$
    \correctchoice $28\pi$ m$^2$
    \choice $40\pi$ m$^2$
\end{choices}

\question[1] Two similar cones have linear scale factor $1:3$. Their volume ratio is:
\begin{choices}
    \choice $1:3$
    \choice $1:6$
    \correctchoice $1:27$
    \choice $1:9$
\end{choices}

\question[1] A ship sails on bearing S\ang{70}W for \SI{40}{km}. Its southward displacement is:
\begin{choices}
    \choice $40\cos\ang{70}$
    \correctchoice $40\cos\ang{70}$ south (approximately \SI{13.7}{km})
    \choice $40\sin\ang{70}$ south
    \choice $40$ km south
\end{choices}

%--- Statistics (4) ---
\question[1] For a normal distribution with $\mu = 80$ and $\sigma = 10$, the $z$-score for
$x = 95$ is:
\begin{choices}
    \choice $0.5$
    \choice $1.0$
    \correctchoice $1.5$
    \choice $2.0$
\end{choices}

\question[1] About 95\% of data in a normal distribution lies within:
\begin{choices}
    \choice $\mu \pm \sigma$
    \correctchoice $\mu \pm 2\sigma$
    \choice $\mu \pm 3\sigma$
    \choice $\mu \pm 0.5\sigma$
\end{choices}

\question[1] A survey of 900 people has a margin of error of $\pm 3.3\%$ at 95\% confidence.
To get a margin of error of $\pm 1.65\%$, the sample size needed is:
\begin{choices}
    \choice $1\,800$
    \choice $2\,700$
    \correctchoice $3\,600$
    \choice $450$
\end{choices}

%--- Applications of Geometry (4) ---
\question[1] The magnitude of vector $\vec{v} = (-5, 12)$ is:
\begin{choices}
    \choice $7$
    \choice $10$
    \correctchoice $13$
    \choice $17$
\end{choices}

\question[1] The perpendicular bisector of segment $AB$ is the locus of points:
\begin{choices}
    \choice On line $AB$
    \correctchoice Equidistant from $A$ and $B$
    \choice Equidistant from $A$ and a fixed line
    \choice Within a fixed distance from $A$
\end{choices}

\question[1] An inscribed angle of \ang{55} subtends an arc. The central angle subtending
the same arc is:
\begin{choices}
    \choice \ang{27.5}
    \correctchoice \ang{110}
    \choice \ang{55}
    \choice \ang{125}
\end{choices}

\question[1] To prove $PQRS$ is a parallelogram using midpoints of diagonals, you verify:
\begin{choices}
    \choice All sides are equal
    \choice The diagonals are perpendicular
    \correctchoice The midpoints of both diagonals are the same point
    \choice All angles are \ang{90}
\end{choices}

\question[1] The shoelace formula gives the:
\begin{choices}
    \choice Perimeter of a polygon
    \correctchoice Area of a polygon given its vertices
    \choice Length of the diagonals
    \choice Number of sides
\end{choices}


\section*{Part B --- Short Answer (40 marks)}

\question[8] \textbf{Financial Mathematics.}
\begin{parts}
    \part[4] Mei contributes \$400 per month to an RRSP at 6\% per year compounded monthly
    for 30 years. Find the total value at retirement and the total amount she contributed.
    How much is interest?
    \[(FV = PMT \cdot \tfrac{(1+i)^n-1}{i})\]
    \part[4] Mei then uses the RRSP to fund a monthly pension. She needs \$2\,500 per month
    for 20 years at 4\% per year compounded monthly. Find the present value required at
    retirement. Does her RRSP cover it?
    \[(PV = PMT \cdot \tfrac{1-(1+i)^{-n}}{i})\]
\end{parts}
\begin{solution}[11cm]
\textbf{(a)} $i = \dfrac{0.06}{12} = 0.005$, $n = 30 \times 12 = 360$
\[FV = 400 \cdot \frac{(1.005)^{360}-1}{0.005} = 400 \cdot \frac{6.0226 - 1}{0.005}
= 400 \cdot \frac{5.0226}{0.005} = 400 \times 1004.52 \approx \mathbf{\$401\,808}\]
Total contributed: $400 \times 360 = \$144\,000$

Interest earned: $\$401\,808 - \$144\,000 = \mathbf{\$257\,808}$

\textbf{(b)} $i = \dfrac{0.04}{12} \approx 0.003\overline{3}$, $n = 20 \times 12 = 240$
\[PV = 2500 \cdot \frac{1-(1.003\overline{3})^{-240}}{0.003\overline{3}}
= 2500 \cdot \frac{1-0.4493}{0.003\overline{3}} = 2500 \cdot \frac{0.5507}{0.003\overline{3}}
\approx 2500 \times 165.02 \approx \mathbf{\$412\,550}\]

Her RRSP of \$401,808 does \textbf{not quite cover} the required \$412,550 — she is short by about \$10,742.
\end{solution}

\question[8] \textbf{Modelling with Functions.}
The monthly average temperature (°C) for a city is given below:

\begin{center}
\begin{tabular}{@{}lrrrrrrrrrrrr@{}}
\toprule
Month & J & F & M & A & M & J & J & A & S & O & N & D \\
\midrule
Temp (°C) & $-8$ & $-6$ & $0$ & 8 & 15 & 20 & 23 & 22 & 16 & 9 & 1 & $-5$ \\
\bottomrule
\end{tabular}
\end{center}
\begin{parts}
    \part[2] Identify the maximum and minimum temperatures. Find the amplitude and
    vertical shift of a sinusoidal model.
    \part[3] The period is 12 months. Find $k$ in $y = a\sin(k(x-d)) + c$. Write a complete
    sinusoidal equation for the temperature, using a sine function with phase shift
    $d = 4$ (April as the starting rise). Use $x = 1$ for January.
    \part[3] Use your model to predict the temperature in month 8 (August) and month 3 (March).
    Compare to the actual values and comment.
\end{parts}
\begin{solution}[10cm]
\textbf{(a)} Max $= 23$°C (July), Min $= -8$°C (January).

Amplitude: $a = \dfrac{23-(-8)}{2} = \dfrac{31}{2} = \mathbf{15.5}$°C

Vertical shift: $c = \dfrac{23+(-8)}{2} = \dfrac{15}{2} = \mathbf{7.5}$°C

\textbf{(b)} $k = \dfrac{2\pi}{12} = \dfrac{\pi}{6}$

Model (using sine with $d = 4$, since the temperature is near the midline rising around April):
\[y = 15.5\sin\!\left(\dfrac{\pi}{6}(x - 4)\right) + 7.5\]

\textbf{(c)} August: $x = 8$:\quad $y = 15.5\sin\!\left(\dfrac{\pi}{6}(4)\right) + 7.5 = 15.5\sin\!\left(\dfrac{2\pi}{3}\right) + 7.5 = 15.5(0.866) + 7.5 \approx \mathbf{20.9}$°C

Actual: 22°C. Difference: 1.1°C — close.

March: $x = 3$:\quad $y = 15.5\sin\!\left(\dfrac{\pi}{6}(-1)\right) + 7.5 = 15.5(-0.5) + 7.5 = -7.75 + 7.5 \approx \mathbf{-0.25}$°C

Actual: 0°C. Difference: 0.25°C — very close.

The sinusoidal model fits the data well. Slight discrepancies occur because real temperature patterns are not perfectly sinusoidal.
\end{solution}

\question[8] \textbf{Trigonometry \& Geometry.}
\begin{parts}
    \part[4] From the roof of a building (height $h$), a surveyor measures the angle of
    elevation to the top of a taller building as \ang{22} and the angle of depression to
    the base of that building as \ang{35}. The horizontal distance between the buildings
    is \SI{80}{m}. Find $h$ (the height of the first building) and $H$ (the total height of
    the second building).
    \part[4] A storage tank is a cylinder with a cone on top. The cylinder has radius \SI{4}{m}
    and height \SI{6}{m}. The cone has the same base radius and slant height \SI{5}{m}.
    Find the total surface area and total volume.
\end{parts}
\begin{solution}[11cm]
\textbf{(a)} Let $d = 80$ m (horizontal distance), $h$ = height of first building.

From angle of depression to the base: $\tan\ang{35} = \dfrac{h}{80}$
\[h = 80\tan\ang{35} = 80(0.7002) \approx \mathbf{\SI{56.0}{m}}\]

From angle of elevation to the top of second building:
Let $x$ = height of second building above the surveyor's level.
$\tan\ang{22} = \dfrac{x}{80} \Rightarrow x = 80\tan\ang{22} = 80(0.4040) \approx \SI{32.3}{m}$

Total height of second building: $H = h + x = 56.0 + 32.3 \approx \mathbf{\SI{88.3}{m}}$

\textbf{(b)} Cone: $r = 4$ m, slant height $l = 5$ m, so cone height $= \sqrt{5^2-4^2} = 3$ m.

\textit{Total Surface Area:}
\begin{itemize}
\item Circular base of cylinder: $\pi(4)^2 = 16\pi$
\item Lateral surface of cylinder: $2\pi(4)(6) = 48\pi$
\item Lateral surface of cone: $\pi r l = \pi(4)(5) = 20\pi$
\item (No circular top on cylinder — the cone sits there)
\end{itemize}
$SA = 16\pi + 48\pi + 20\pi = 84\pi \approx \mathbf{263.9 \text{ m}^2}$

\textit{Total Volume:}
$V_{\text{cyl}} = \pi(4)^2(6) = 96\pi$;\quad $V_{\text{cone}} = \tfrac{1}{3}\pi(4)^2(3) = 16\pi$

$V = 96\pi + 16\pi = 112\pi \approx \mathbf{351.9 \text{ m}^3}$
\end{solution}

\question[8] \textbf{Statistics.}
The lifespans (years) of a species of parrot are normally distributed with $\mu = 35$ and $\sigma = 6$.
\begin{parts}
    \part[3] Find the percentage of parrots that live:
    (i) between 29 and 41 years;\quad
    (ii) more than 47 years;\quad
    (iii) less than 23 years.
    \part[3] Find the lifespan at the 75th percentile ($z_{0.75} = 0.674$) and the lifespan
    below which 10\% of parrots fall ($z_{0.10} = -1.282$).
    \part[2] A wildlife researcher surveys 100 parrots and finds a mean lifespan of 36.2 years.
    The margin of error at 95\% confidence is $\pm \dfrac{1.96 \times 6}{\sqrt{100}}$.
    Find the confidence interval and interpret it.
\end{parts}
\begin{solution}[11cm]
\textbf{(a)}
(i) $29 = 35 - 6 = \mu - \sigma$ and $41 = 35 + 6 = \mu + \sigma$. By the 68-95-99.7 rule: $\mathbf{68\%}$

(ii) $47 = 35 + 2(6) = \mu + 2\sigma$. Above $\mu + 2\sigma$: $\dfrac{100-95}{2} = \mathbf{2.5\%}$

(iii) $23 = 35 - 2(6) = \mu - 2\sigma$. Below $\mu - 2\sigma$: $\mathbf{2.5\%}$

\textbf{(b)} 75th percentile: $x = \mu + z\sigma = 35 + 0.674(6) = 35 + 4.04 = \mathbf{39.0}$ years

10th percentile: $x = 35 + (-1.282)(6) = 35 - 7.69 = \mathbf{27.3}$ years

\textbf{(c)} $ME = \dfrac{1.96 \times 6}{\sqrt{100}} = \dfrac{11.76}{10} = 1.176$

Confidence interval: $36.2 \pm 1.176 \Rightarrow \mathbf{(35.02,\; 37.38)}$ years

Interpretation: We are 95\% confident that the true mean lifespan of all parrots of this species is between 35.0 and 37.4 years.
\end{solution}

\question[8] \textbf{Applications of Geometry.}
\begin{parts}
    \part[4] A ship leaves port $O$ and sails on bearing N\ang{50}E for \SI{120}{km} to point $A$.
    It then sails on bearing S\ang{30}E for \SI{90}{km} to point $B$.
    Find the distance $OB$ and the bearing from $O$ to $B$.
    \part[4] The vertices of quadrilateral $PQRS$ are $P(0,1)$, $Q(4,3)$, $R(6,0)$, $S(2,-2)$.
    Use midpoints of diagonals to prove $PQRS$ is a parallelogram. Then find its area using the
    shoelace formula.
\end{parts}
\begin{solution}[11cm]
\textbf{(a)} The angle at $A$ between the two legs:

Bearing N\ang{50}E (first leg) and S\ang{30}E (second leg):

The first leg arrives at $A$ from the direction S\ang{50}W.
The second leg departs on S\ang{30}E.
The angle between these two directions at $A$:
$\angle OAB = 180^\circ - 50^\circ - 30^\circ = \ang{100}$

By Cosine Law:
\[OB^2 = 120^2 + 90^2 - 2(120)(90)\cos\ang{100}\]
\[= 14\,400 + 8\,100 - 21\,600(-0.1736) = 22\,500 + 3\,750 = 26\,250\]
\[OB = \sqrt{26\,250} \approx \mathbf{\SI{162.0}{km}}\]

For bearing: $\dfrac{\sin(\angle AOB)}{90} = \dfrac{\sin\ang{100}}{162.0}$

$\sin(\angle AOB) = \dfrac{90 \times 0.9848}{162.0} = 0.5463 \Rightarrow \angle AOB \approx \ang{33.1}$

Bearing from $O$ to $B$: N\ang{50}E rotated \ang{33.1} clockwise $=$ N\ang{83}E (approximately).

\textbf{(b)} Midpoint of $PR$: $\left(\dfrac{0+6}{2}, \dfrac{1+0}{2}\right) = (3, 0.5)$

Midpoint of $QS$: $\left(\dfrac{4+2}{2}, \dfrac{3+(-2)}{2}\right) = (3, 0.5)$ \checkmark

Diagonals bisect each other, so $PQRS$ is a parallelogram.

Shoelace formula:
\[
\text{Area} = \tfrac{1}{2}|x_P(y_Q - y_S) + x_Q(y_R - y_P) + x_R(y_S - y_Q) + x_S(y_P - y_R)|
\]
\[= \tfrac{1}{2}|0(3-(-2)) + 4(0-1) + 6((-2)-3) + 2(1-0)|\]
\[= \tfrac{1}{2}|0 + (-4) + (-30) + 2| = \tfrac{1}{2}|-32| = \mathbf{16 \text{ square units}}\]
\end{solution}

\end{questions}
\end{document}
